\subsection*{A. Reproduksi}
\begin{verbatim}
pip install -r kode/requirements.txt
python -m pytest -q kode/tests
python kode/run_all.py --task all --mode train
python kode/run_all.py --task t1 --mode reproduce
\end{verbatim}
Perintah ketiga menghasilkan semua angka, gambar, tabel, dan tiga file prediksi (GPU T4: ±2--2,5 jam); perintah keempat memuat model tersimpan dan menghitung ulang
angka kunci dalam beberapa menit. Seed 42; versi library dan GPU tercatat di setiap \url{kode/outputs/results/*.json}.

\subsection*{B. Penggunaan AI}
Kode disusun dengan bantuan Claude (Anthropic). Kelompok memverifikasinya dengan tes otomatis (metrik, anti-kebocoran jendela RUL), sanity check overfit satu batch,
membandingkan skema validasi, menjalankan ulang \texttt{-{}-mode reproduce}, dan memeriksa setiap tabel serta gambar. Rincian per anggota ada di \texttt{kontribusi.md}.
